%% article.tex -- draft for submission to Ecological Modelling.
%%
%% STATUS: section plan with worked figures. Prose marked [OUTLINE] is a plan of
%% what the section must argue, not finished text.
%%
%% Submission format: Elsevier wants elsarticle.cls. That class is not in the
%% build container, so this drafts in `article` and switches at submission:
%%   \documentclass[preprint,12pt]{elsarticle}
%% with \begin{frontmatter}...\end{frontmatter} replacing \maketitle.
%%
%% Figures come from the fidelity harness, not from hand-drawn illustrations:
%%   make compare sheet     # regenerates every figure this file includes
%% That is deliberate. The figures showing that the symbols match Odum's
%% originals are produced by the same code that tests it.

\documentclass[11pt,a4paper]{article}
\usepackage[margin=1in]{geometry}
\usepackage{graphicx}
\usepackage{subcaption}
\usepackage{booktabs}
\usepackage{listings}
\usepackage{natbib}
\usepackage{energese}
\usepackage{hyperref}

% Both spellings, so the paper builds whether lualatex runs from the repository
% root (as the Makefile does) or from docs/ (as an editor's build button does).
\graphicspath{{tests/fidelity/out/}{reference-images/}%
              {../tests/fidelity/out/}{../reference-images/}}

\lstset{basicstyle=\ttfamily\footnotesize, columns=fullflexible,
        breaklines=true, frame=single, framesep=4pt, xleftmargin=6pt}

\title{Energese: reproducible typesetting of Odum Energy Systems Language
       diagrams in \LaTeX}
\author{S. Maud}
\date{Draft}

\begin{document}
\maketitle

\begin{abstract}
[OUTLINE] State the problem: Energy Systems Language diagrams are drawn by hand
in general-purpose vector tools, so they are not reproducible, not
version-controllable, not generated from the models they depict, and drift from
Odum's published symbol conventions. State the contribution: \texttt{energese},
a Lua\LaTeX{} package that renders ESL diagrams from a declarative model, in
either JSON or \TeX, through a single layout engine; and a two-tier fidelity
method that makes conformance to Odum's symbols a machine-checked property
rather than a matter of opinion. State the result quantitatively: all nine
implemented symbols agree with geometry measured from Odum's published sheet to
within measurement precision on two independent invariants, and the two
authoring front ends are verified to produce byte-identical raster output for
models up to and including a complete published figure. State the availability: CTAN, with the
model format shared with the GSSK simulation kernel so one file both simulates
and draws.
\end{abstract}

\noindent\textbf{Keywords:} energy systems language; emergy; systems ecology;
scientific typesetting; reproducible figures; H.T. Odum

\section{Introduction}
\label{sec:intro}

[OUTLINE] Three moves.

\emph{Why ESL matters.} Odum's Energy Systems Language is a formal notation, not
an illustrative convention: each symbol carries defined semantics about energy
quality, storage, interaction and dissipation \citep{odum1994,odum1996}. A
diagram is a specification.

\emph{Why current practice undermines that.} In practice ESL diagrams are drawn
by hand. Consequences to argue, each of which this package addresses: symbol
geometry drifts from the published convention with no mechanism to detect it;
diagrams cannot be regenerated when the underlying model changes, so figure and
model silently diverge; a diagram is a binary artifact, so it cannot be
meaningfully diffed or reviewed; and the diagram carries no machine-readable
model, so the reader cannot recover what was drawn.

\emph{What this paper contributes.} (i) A package that generates ESL diagrams
from a declarative model. (ii) Dual front ends over one engine, so a model can
live in a data file or in the document. (iii) A fidelity method that separates
what can be checked exactly from what can only be checked within measurement
tolerance, and does both automatically. (iv) A shared model format with a
simulation kernel.

\section{Background: the Energy Systems Language}
\label{sec:background}

\subsection{Origins and purpose}

[OUTLINE] Odum developed ESL from the 1950s onward, grounding it deliberately in
electrical circuit schematics \citep{odum1971,odum1994}. The grounding is the
point: a circuit schematic is not decoration but a statement about conserved
quantities, and the analogy was meant to carry that discipline into ecology.
Each symbol denotes a class of process with defined behaviour --- storage,
production, consumption, interaction, exchange --- and each pathway a flow.

Say why he thought it necessary. Odum's recurring complaint was that scholars
were building mathematical models of systems which were energetically invalid:
formally well-posed, numerically tractable, and impossible. A notation whose
symbols carry thermodynamic commitments was intended to make such models harder
to write down.

\subsection{Prior tooling, and the distribution problem}
\label{sec:tooling}

[OUTLINE] Trace why the notation never acquired reliable tooling. ESL symbol
sets have circulated as stencils and templates for general-purpose drawing
programs --- Microsoft Visio, CorelDRAW, Adobe Illustrator --- which are
proprietary, individually licensed, and store diagrams in closed formats. Cost
and format both bear on who can participate and on whether a published figure
can be regenerated. More recently draw.io offered an open, browser-based
alternative.

Report the author's own attempt, because its failures motivated this work: an
ESL template built in draw.io at the request of Charles A.\ Hall for his
doctoral students.\footnote{% [PLACEHOLDER] Repository URL for the draw.io ESL
% template, and the form of acknowledgement Hall would prefer.
Repository reference to be added.} Two things went wrong, and neither was
technical. Distribution reached one research group and stopped, because a
template is a file rather than a package with a registry behind it. And the
template could validate nothing: it drew whatever was dragged onto it.


Draw the general conclusion: a symbol library needs no semantics, so nobody
building one is forced to supply any. That is why the recurring difficulty has
been maintaining and distributing a standard symbol set, and why the notation's
implied rules have stayed implicit for fifty years.

\section{Related work}
\label{sec:related}

[OUTLINE] Tell the tooling history, because it explains why the problem has
persisted. ESL symbol sets have been distributed as stencils and templates for
proprietary drawing tools --- Microsoft Visio, CorelDRAW, Adobe Illustrator ---
whose cost and closed formats limit both access and reproducibility. More
recently draw.io offered an open, web-based alternative, and the present author
built an ESL template in it at the request of Charles A. Hall for his doctoral
students. That exercise motivated this package, and its two failures are
instructive: distribution was limited to one research group, and the template
could not validate anything it drew. Maintaining and distributing a standard
symbol set has proven the recurring difficulty.

Position against: general diagramming tools (Inkscape,
\textsc{OmniGraffle}) --- expressive, not reproducible, no domain semantics;
\TeX{} drawing packages (TikZ, PSTricks) --- reproducible but with no ESL
vocabulary, so every author re-derives the symbols; systems-dynamics tools
(Stella, Vensim, Insight Maker) --- these do carry model semantics and can draw,
but in their own notation rather than Odum's, and their output is not
typeset-quality nor \LaTeX-native; and existing emergy tools, which mostly
compute rather than draw. The gap: nothing renders \emph{Odum's} notation from a
model, reproducibly, inside a document.

\section{ESL as an informal language with convention-based syntax and semantics}
\label{sec:informal}

[OUTLINE] This section carries the paper's central claim, and it should come
before the implementation rather than after it.

\subsection{A semiotic frame}
\label{sec:semiotics}

[OUTLINE] Borrow Morris's trichotomy \citep{morris1938} to say precisely what
ESL has and lacks, because ``informal'' on its own is too blunt.

\emph{Syntactics} --- the relations among signs. ESL has strong conventions
here: transformity orders the horizontal axis, dissipation runs downward,
feedback arcs over the top. None is written down as a rule.

\emph{Semantics} --- the relation of signs to what they denote. This is where
ESL is strongest: each symbol has a defined process meaning, and Odum is
explicit about it.

\emph{Pragmatics} --- the relation of signs to their users. Also strong, and
under-acknowledged: ESL is drawn to be argued over in seminars, and its
conventions serve legibility as much as denotation.

The gap is therefore specifically syntactic. Goodman's conditions on notational
systems \citep{goodman1976} sharpen the diagnosis: a notation requires syntactic
disjointness and finite differentiation --- one must be able to decide whether
two inscriptions are instances of the same character. ESL fails this for
\emph{position}, which carries meaning but has no defined scale, while passing
comfortably for its symbols. Moody's evaluation of visual notations
\citep{moody2009} supplies the vocabulary for the rest, and Harel's argument
that visual formalisms can be as rigorous as textual ones \citep{harel1988}
establishes that the gap is contingent, not inherent. Bertin
\citep{bertin1983} supplies the account of position as a visual variable.

Frame the contribution accordingly: we are not proposing that ESL be replaced
by a formal language, but that its existing conventions be stated precisely
enough to be implemented and checked --- and we report what happens when one
tries.

ESL is a graphical language but not a formal one. It has no grammar and no
stated semantics, so a diagram cannot be checked, and a tool cannot be written
against a specification --- only against custom. That is why every previous
attempt has been a symbol library for a general-purpose drawing tool
(\S\ref{sec:related}) rather than a language implementation: a symbol library
needs no semantics.

Building a renderer removes that option. Placing a node requires deciding what
position \emph{means}; validating one requires deciding what a well-formed
diagram \emph{is}. The decisions below were forced by implementation, and
stating them is the contribution --- not because they are the only possible
answers, but because they are currently unstated anywhere.

\subsection{What the implied syntax asserts, and what it leaves open}

[OUTLINE] Transformity orders the horizontal axis: low quality left, high
quality right. That much is consistently observed in Odum's figures and stated
informally in the literature. What is \emph{not} stated is the scale. Placement
is plainly not linear --- transformity spans orders of magnitude --- and a
logarithmic axis would assert a metric the language never defines.

Argue for the answer this package implements: placement is \emph{ordinal}.
Distinct transformity values are ranked and mapped to evenly spaced columns.
This encodes exactly what the syntax claims, the ordering, and nothing further.
Nodes of equal transformity share a column, which makes a column an equivalence
class of energy quality.

The vertical axis is the harder question. Consider first, and reject, the
obvious candidate: that height encodes power, the capacity to do work, with
dissipated heat low and forcing functions high.

The objection is dimensional. Transformity is emergy per unit available energy,
so along a transformation chain, where emergy is carried forward and available
energy is degraded, transformity rises precisely \emph{because} power falls;
empower is approximately conserved. Available power is therefore already
determined --- inversely --- by horizontal position. Encoding it again on the
vertical axis would be redundant where the two agree and contradictory where
they do not, which is exactly the double-encoding Bertin warns against
\citep{bertin1983}. Odum's own figures confirm it: the Sun, by far the largest
power input to an ecosystem, is drawn at the \emph{bottom} of the source column,
not the top.

The reading that survives is simpler, and unifies the two axes. The vertical
carries the same quality ordering as the horizontal, applied \emph{relative to
the main chain} rather than absolutely: a component of higher transformity than
the chain node it connects to sits above it --- a high-quality input, a
controller, or a higher-quality product --- and degraded output sits below. The
energy hierarchy therefore runs from bottom-left to top-right rather than simply
left to right, and the Sun's position at the bottom of the source column is the
same rule operating where the horizontal axis cannot separate components that
all lie outside the frame.

Report the two tests, because the claim is falsifiable and was tested rather
than asserted. Stripping every hand-placed coordinate from the aggregated
economy, the layout engine raises the high-transformity fuel source above the
chain unprompted, reproducing Odum's placement of Fuels and Minerals entering
from the top. Stripping both coordinates and hand-written layer annotations from
a second model, the rule recovers the author's \texttt{decomposition}
annotation on a waste store from transformity alone.

Note one refinement the tests forced: the chain is built from energy pathways
only. Counting material pathways let a waste store capture the chain and sit on
the axis, when a flow of degraded material to a decomposition store is precisely
a side branch --- the same distinction that exempts material flows from the
quality check in \S\ref{sec:validity}, arrived at independently.

Note the apparent anomaly and dissolve it: the heat sink sits at bottom centre,
which looks like a violation of the horizontal ordering. It is not, because the
ground carries no transformity at all --- it is a terminator, not a component.
Dissipation going down is the vertical convention, not an exception to the
horizontal one.

\subsection{A worked anomaly: the ground symbol}
\label{sec:ground}

[OUTLINE] Use the package to settle a question about the notation, which is the
method this paper is arguing for: a diagram whose rules are implemented can be
interrogated by drawing.

The apparent anomaly: if the horizontal axis orders components by transformity,
the heat sink at bottom centre is a violation, because dissipated energy is the
lowest quality in the system and belongs at the far left.

The resolution: the ground is not a component. It carries no transformity, takes
no part in the quality ordering, and terminates pathways rather than
transforming anything. Figure~\ref{fig:ground} shows the same three components
with dissipation drawn, and the ground sitting outside the ordering entirely.
Dissipation running downward is the vertical convention doing its work; the
horizontal axis is untouched.

\begin{figure}[htbp]
\centering
\begin{energese}[scale=0.85]
    \esmeta{column_spacing=3.2, row_spacing=1.6}
    \esnode{sun}{source}     {Tr=1,    label={Sun}}
    \esnode{grass}{producer} {Tr=100,  label={Grass}}
    \esnode{cow}{consumer}   {Tr=2000, label={Cow}}
    \esflow{sun}{grass}{energy}
    \esflow{grass}{cow}{energy}
\end{energese}
\caption{Three components ordered left to right by transformity, with their
dissipation pathways converging on a single ground. The ground carries no
transformity and is placed by neither rule: it terminates the diagram rather
than occupying a position in it. The pathways converge as a delta and not as a
spray: how far each drops before turning in is ordered by its distance from the
trunk, so the furthest runs below and outside the nearest all the way to the
junction they share. Two dissipation pathways that cross read as pathways
meeting where nothing joins.}
\label{fig:ground}
\end{figure}

Generalise the move: three node types --- \texttt{ground}, \texttt{text} and
\texttt{box} --- are annotations or terminators rather than components, and are
exempt from carrying transformity. Making that exemption explicit is what
dissolves the anomaly, and it is the kind of distinction an implementation
forces and a drawing tool never would.

\subsection{Validity, and where it stops}
\label{sec:validity}

[OUTLINE] Odum's stated motivation for ESL was that scholars build mathematical
models which are energetically invalid; the language was grounded in circuit
schematics to help them avoid it. Drawing tools cannot help, because they render
whatever is drawn.

Set out what a renderer \emph{can} check without leaving the diagram: quality
does not fall along an energy pathway; sources take no inflow and grounds emit
none; every component is reachable from a source; energy that arrives leaves,
as outflow or as dissipation. These are topological, need no quantities, and
correspond directly to Odum's complaint.

Then set out the boundary, which matters as much as the checks:

\begin{itemize}
  \item \emph{Material pathways are exempt from the quality rule.} Degraded
        material moving to a waste or decomposition store genuinely falls in
        transformity --- that is what decomposition is. Applying the rule to
        material flows produced a false positive on one of the package's own
        test models, which is worth reporting as evidence that the distinction
        is real rather than stipulated.
  \item \emph{Conservation is not checked.} Balancing inflow against outflow
        needs quantities, and a diagram carries none. That is the seam with a
        simulation kernel (\S\ref{sec:gssk}).
  \item \emph{Emergy is out of scope entirely.} The algorithms are not
        standardised and there is no formal definition to validate against, so
        a check would assert a standard that does not exist. The package
        represents the pre-emergy form of these models; the dual energy/emergy
        depiction of Valyi's emergy simulator is not supported.
\end{itemize}

Report the finding honestly: the quality check caught a genuine error in this
paper's own worked example, where an interaction fed by a transformity-1000
source had been given a transformity of 500.

\section{Design}
\label{sec:design}

\subsection{One model, two front ends}

The central design constraint is that there is exactly one model and one
renderer. A system is described by nodes, edges and metadata. That description
may be written as JSON, for models that live outside the document or are
generated by a simulation, or directly in \TeX, for diagrams that belong with
their prose. Both are parsed into the same Lua table and rendered by the same
code.

Figure~\ref{fig:frontends} shows the same system expressed both ways.

\begin{figure}[htbp]
\begin{minipage}[t]{0.48\textwidth}
\begin{lstlisting}[title=JSON]
{"metadata": {"column_spacing": 3.0},
 "nodes": [
  {"id":"sun",   "type":"source",
   "Tr":1,    "label":"Sun"},
  {"id":"grass", "type":"producer",
   "Tr":100,  "label":"Grass"},
  {"id":"cow",   "type":"consumer",
   "Tr":2000, "label":"Cow"}],
 "edges": [
  {"from":"sun","to":"grass",
   "type":"energy"},
  {"from":"grass","to":"cow",
   "type":"energy"}]}
\end{lstlisting}
\end{minipage}\hfill
\begin{minipage}[t]{0.48\textwidth}
\begin{lstlisting}[title=\TeX]
\begin{energese}
  \esmeta{column_spacing=3.0}

  \esnode{sun}{source}
         {Tr=1,    label={Sun}}
  \esnode{grass}{producer}
         {Tr=100,  label={Grass}}
  \esnode{cow}{consumer}
         {Tr=2000, label={Cow}}

  \esflow{sun}{grass}{energy}
  \esflow{grass}{cow}{energy}
\end{energese}
\end{lstlisting}
\end{minipage}

\vspace{1em}
\centering
\begin{energese}[scale=0.8]
  \esmeta{column_spacing=3.0, row_spacing=1.5}
  \esnode{sun}{source}     {Tr=1,    label={Sun}}
  \esnode{grass}{producer} {Tr=100,  label={Grass}}
  \esnode{cow}{consumer}   {Tr=2000, label={Cow}}
  \esflow{sun}{grass}{energy}
  \esflow{grass}{cow}{energy}
\end{energese}

\caption{The same system in both front ends, and the diagram they produce. The
figure is rendered by the package as this paper is typeset. Equivalent models
are verified to rasterise identically (\S\ref{sec:parity}).}
\label{fig:frontends}
\end{figure}

\subsection{Routing}

[OUTLINE] Make the semantic point first, because it is easy to mistake for a
cosmetic one. In ESL a pathway meeting a symbol denotes flow entering that
symbol. A pathway drawn across an unrelated symbol therefore asserts a
connection the model does not contain --- it is a false statement, not an
untidy one. Drawing the pathway on a lower layer so it passes behind the symbol
does not help: it still reads as flow into the symbol, and merely hides the
evidence. The pathway has to move.

Describe the router: a pathway takes the shortest route that avoids every
symbol which is not its own endpoint, found as a shortest path through a
visibility graph over the corners of the obstacles, inflated by the clearance a
pathway keeps. For rectangular obstacles that path is \emph{the} shortest one
available, since a taut string around boxes bends only at their corners.
Waypoints are joined by curves whose tangents follow the chord through each
point, so the pathway is continuous in tangent at every waypoint and flows
through a detour rather than turning a corner at it.

Worth saying why the obvious cheaper thing does not work, because it is what
this replaced: stepping over one offending symbol at a time, to whichever side
deviates less, is greedy in the ordinary sense and wrong in the ordinary way ---
clearing the first obstacle by the smaller deviation regularly pushes the
pathway into a second, and what comes out of a crowded middle is a staircase of
local decisions with no claim to being short.

One further distinction that costs a paragraph and is invisible until it is
wrong: what must clear the symbol is the \emph{curve}, not the polyline the
router returns. The drawn curve cuts the corners the polyline turns at, so a
polyline that clears an obstacle by a hair produces a pathway that crosses it.
The router therefore tests the curve it is about to hand to Ti\emph{k}Z, and
widens the corners it bends around until the curve clears too.

Note the limitation honestly: obstacle extents are estimated from the grid and a
per-character label estimate, because the layout engine runs in Lua and cannot
know what \TeX{} will typeset. The estimate is not merely the label's width ---
each shape declares what fraction of its outline the text may occupy, and a
labelled symbol is wider than its label by exactly that factor. Exact extents
would still require round-tripping node dimensions through the \texttt{.aux}
file.

\subsection{Labels on pathways}

[OUTLINE] A flow's quantity is content, and a quantity written across a symbol,
across another pathway, or across the system frame is content the reader cannot
recover --- worse when it lands between two pathways and they have to guess
which it belongs to. Each label is placed by scoring the positions Ti\emph{k}Z
can express for a node on a path --- a fraction along it, above or below, at one
of a few distances clear --- against the symbols, the pathways sampled as the
curves they will be drawn as, the dissipation bundle, the frame, and the labels
already placed.

Argue for enumeration over the annealing the graph-drawing literature reaches
for. Two reasons, and the second is the stronger. Reproducibility: a golden
image is only a test if the same input renders identically, and a randomised
search would have to carry a seed and would still move every label whenever
anything upstream of it changed. And expressibility: the candidate set is not a
discretisation of a continuous space, it \emph{is} the set of placements
Ti\emph{k}Z can express relative to a path. A position it cannot express has to
be given absolutely, and an absolutely-placed label stops travelling with its
pathway the moment the layout moves.

Report the residue honestly. Where two symbols stand closer together than the
label is wide, no candidate is clear and the search returns the least bad one;
Figure~\ref{fig:complex-json} has such a case, where a quantity sits on the
system frame because every alternative sits on a symbol. The placement is also
approximate in the same way the extents are: Lua scores the label against the
pathway between two node boundaries, while \TeX{} draws it against the pathway
between two ports.

\subsection{Layout}

[OUTLINE] Explain the two axes and why they are chosen this way.

Horizontal position derives from transformity: distinct $Tr$ values are sorted
and assigned to columns, so horizontal position expresses energy quality rather
than declaration order --- which is the convention readers of ESL already apply.
Vertical position uses a central spine: the longest chain of energy and material
pathways is laid on the axis, and remaining nodes are distributed above and
below their column's spine. Explain why automatic layout matters here: it is
what makes a diagram a \emph{function} of its model rather than an artifact
maintained alongside it. Note the escape hatch --- explicit coordinates override
either axis independently, so one awkward component does not force abandoning
the layout.

Then the exception, which is a claim about the notation rather than about
layout. An interaction is not an independent component: it is where two flows
combine on their way into the thing they drive, and Odum draws it against that
thing. Given a column of its own by its transformity, it drifts away from the
component it modifies and the reader loses which one that is. Interactions and
standalone transactions are therefore children in a compound graph --- locked to
a parent by a fixed offset, and exempt from the column algorithm, so they
neither claim a column nor push a sibling out of a row. The offset follows the
flow: a child that feeds its parent sits upstream of it, one fed by its parent
sits downstream, because a modifier drawn on the wrong side reverses the
direction its pathway reads in. The gap is half a column or whatever the two
symbols need to stand clear of each other, whichever is larger --- a labelled
exchange diamond is wide, and a fixed fraction of the column buries its left
vertex inside the consumer it belongs to.

\subsection{The system window}

[OUTLINE] The frame declares what the model treats as inside the system, so
what it encloses is a claim and not a matter of taste. Two rules settle it.
What is inside: everything but the sources, which is Odum's convention --- a
source is by definition an input from outside --- overridable per node. And
where the frame falls: a fixed margin outside the \emph{extents} of what is
inside, labels included, rather than outside their centres. Measuring centres
and padding by half a column, as an earlier version did, makes the margin a
function of the widest symbol in the diagram: the frame crowds a wide producer
and leaves a hand's width of dead space beside a small storage.

The frame must also leave room underneath for the dissipation bundle, which is
the part that is easy to miss. The funnel converges \emph{inside} the window and
only its trunk crosses the frame --- what leaves the system is one aggregate
flow of used energy, not several, and letting each pathway cross separately both
shreds the boundary and misstates the accounting. A frame fitted to the symbols
alone cuts straight across that funnel.

\subsection{Symbols that size to their content, on a common grid}

[OUTLINE] Two rules, and the second is the one that took longest to get right.

\emph{A symbol sizes itself to its label.} Fixed-size symbols are why hand-drawn
ESL diagrams have labels spilling outside their shapes or shrunk to
illegibility.

\emph{A symbol keeps its shape while doing so.} Growing to fit a label must be
uniform: a hexagon widened to swallow a long caption reads as a different
symbol, and the notation's whole claim is that shape carries meaning. Each shape
therefore declares its measured aspect ratio and the fraction of its outline the
text may occupy, and the outline is the smallest size at that ratio which
contains the label. Where a label will not fit even then, the label is displaced
rather than the shape distorted. Report the check: consumers carrying a single
character, one word, a two-line caption, and a magnitude of 1.5 all measure
between 0.874 and 0.880 against a canonical 0.877.

Note the intermediate version that was wrong, because it is instructive: capping
how far a label could flatten a symbol merely bounded the distortion instead of
removing it, and the resulting hexagons still read as squashed. Each symbol here
derives its geometry from its own label and grows the outline around it.

Then the complication, which is the interesting part. Each shape grows from its
label at a \emph{different} rate, because each encloses a rectangle differently:
a producer's height is twice the label's half-height, a consumer's is $3.23$
times, an amplifier's $5.2$ times. Sizing each symbol independently therefore
makes identical labels produce wildly different symbols --- which looks like a
rendering fault but is a direct consequence of the geometry.

The resolution is a grid. A single unit, \verb|\energeseunit|, sets a floor for
every symbol, and each symbol's floor is the multiple of that unit at which its
bare form lands on the relative height Odum draws it at
(Table~\ref{tab:conformance}). Short labels fall inside the floor, so every
symbol sits on the grid; a label too large for its cell grows its own symbol,
which is the only way to keep the text readable. The grid is never drawn.

State that this is asserted, not asserted-to: rendered heights agree with Odum's
measured ratios to within $0.02$.

\section{Method: making fidelity checkable}
\label{sec:method}

This is the methodological contribution and the section to write most carefully.

\subsection{Why ``pixel-perfect'' needs defining}

[OUTLINE] The obvious formulation --- compare our output to Odum's published
figures pixel by pixel --- is not a well-posed test, and saying so precisely is
part of the contribution. The available references are $144$\,dpi captures of
printed figures in which individual symbols span $31$--$107$ pixels. At that
resolution a linear dimension is recoverable to roughly $1$--$3\%$; stroke
weight, curve control points and joins are not recoverable at all. The captures
also carry scanner artifacts and, in one case, visible show-through from the
reverse of the page. Any threshold that ``passed'' such a comparison would be
measuring noise. Cite this as a general problem for reproducing figures from the
printed literature.

\subsection{Two tiers}

We therefore separate two questions that are usually conflated.

\emph{Tier 1, regression:} has the output changed? The reference is the
package's own approved output. Rendering is deterministic --- verified: rasters
are byte-identical between runs, and with \texttt{SOURCE\_DATE\_EPOCH} pinned
the PDFs are too --- so any difference at all is a real change, and comparison
is exact. A changed render fails the build and writes a three-panel diff.
Accepting a change is explicit and versions the new reference alongside the code.

\emph{Tier 2, conformance:} is the output right? Here the reference is Odum's
published symbol sheet, and the assertions are scale-independent geometric
invariants measured from it, with tolerances no tighter than the measurement
precision justifies.

Tier 2 establishes that a symbol is correct; Tier 1 ensures it stays correct.

\subsection{Measuring the reference}

[OUTLINE] Describe the extraction and why the naive approach fails. Odum draws
each symbol with inflow, outflow and heat-sink arrows attached, so a bounding
box around the ink measures the arrows too. We instead flood-fill each symbol's
interior from a seed point: the outline is closed and the arrows lie outside it,
so the filled region is the symbol body alone. Seed points are the only manual
input; everything downstream is measured and re-derivable.

State the internal validation: the source symbol is a circle, and the procedure
recovers an aspect ratio of $0.984$ against a true value of $1$. That $1.6\%$
error is an empirical estimate of the method's precision, and it sets the floor
for every tolerance in Table~\ref{tab:conformance}.

\subsection{Results}

\begin{table}[htbp]
\centering
\begin{tabular}{lrrrrrr}
\toprule
& \multicolumn{2}{c}{Aspect ratio} & \multicolumn{2}{c}{Interior overlap}
& \multicolumn{2}{c}{Height / producer} \\
\cmidrule(lr){2-3}\cmidrule(lr){4-5}\cmidrule(lr){6-7}
Symbol & Ref. & Rend. & IoU & Min. & Ref. & Rend. \\
\midrule
source        & 0.984 & 1.000 & 0.982 & 0.93 & 1.016 & 1.015 \\
storage       & 0.968 & 1.007 & 0.901 & 0.85 & 1.000 & 1.000 \\
producer      & 1.726 & 1.701 & 0.990 & 0.94 & 1.000 & 1.000 \\
consumer      & 0.877 & 0.879 & 0.966 & 0.92 & 1.306 & 1.299 \\
transaction   & 2.226 & 2.130 & 0.669 & 0.63 & 0.500 & 0.515 \\
interaction   & 1.886 & 1.844 & 0.934 & 0.89 & 0.565 & 0.575 \\
gain          & 0.917 & 0.925 & 0.919 & 0.87 & 0.581 & 0.597 \\
loop\_limited & 1.019 & 1.017 & 0.962 & 0.91 & 0.855 & 0.858 \\
switch        & 1.036 & 1.033 & 0.922 & 0.87 & 0.903 & 0.910 \\
\bottomrule
\end{tabular}
\caption{Three independent invariants for each symbol. Aspect ratio is width over
height of the symbol's interior. Interior overlap is intersection-over-union of
the enclosed regions after normalisation --- interiors rather than outlines, so
the measure is independent of stroke weight. Thresholds are the measured
agreement less $0.05$, so each reflects what that symbol pair actually achieves.
Height is relative to the producer, and is what the grid guarantees. The
exchange diamond is the weakest: at $31$ pixels tall the reference's own stroke
consumes much of its interior, which caps the achievable overlap.}
\label{tab:conformance}
\end{table}

\begin{figure}[htbp]
\centering
\includegraphics[width=\textwidth]{compare-all.png}
\caption{Odum's published symbol (left of each pair) beside the package's
rendering (right). The reference panels are cropped by the same flood-fill that
produced the measurements in Table~\ref{tab:conformance}, so the figure shows
exactly what was measured. Each pair is scaled to a common height so that shape
can be compared; stroke weight therefore cannot be, since a symbol spanning
$31$ pixels in the source is magnified far more than one spanning $107$.}
\label{fig:sidebyside}
\end{figure}

\subsection{A worked case: the method finding a defect}
\label{sec:defect}

[OUTLINE] This is the strongest evidence that the method does work, and it
should be reported honestly rather than smoothed over.

During development the consumer symbol was implemented as a hexagon pointed left
and right, from a reading of the full symbol sheet at its native resolution. The
existing test suite --- nine compile-and-check suites --- passed throughout,
because compiling successfully says nothing about shape. The conformance
measurement put the rendered symbol at an aspect ratio of $1.435$ against a
reference of $0.877$: the symbol is in fact pointed top and bottom with vertical
sides, through which energy enters and leaves. Correcting the topology brought
the rendering to $0.876$.

Points to draw out: a compile-only test suite cannot detect a wrong symbol; the
error survived visual inspection of the rendered diagrams because a hexagon
looks plausible either way; and the defect was found by measurement, not by
looking harder. Generalise this to reproducible-figure practice.

\subsection{Front-end equivalence}
\label{sec:parity}

[OUTLINE] Describe the property and how it is enforced. Because both front ends
are claimed to be interchangeable, the claim is tested rather than asserted:
equivalent models are rendered through each path, rasterised and compared
exactly. Report that this check found a real defect --- the JSON path was
emitting stray spaces that padded every diagram's bounding box by
approximately $36$\,pt, invisible in isolation and detectable only by comparison.

\subsection{A complete model in both front ends}
\label{sec:complex}

Figure~\ref{fig:frontends} is deliberately small. Figure~\ref{fig:complex-tex}
gives a complete published diagram --- Odum's aggregated economy --- authored in
\TeX, and Figure~\ref{fig:complex-json} the same model as JSON. Both listings are
included directly from the repository's example files, so they cannot drift from
what is tested; and \texttt{make parity} asserts that the two render identically.

Between them they exercise every feature the package has: manual coordinates
where transformity alone will not place a node, a system window, a relocated
heat sink, interaction symbols, bent pathways, and several money-feedback arcs.

Two modelling choices in this figure are worth stating, because neither is
forced by the source and a reader comparing the two will notice.

\emph{Interaction symbols are made explicit.} Odum draws a notched shape abutting
the left of each consumer, where separate inflows must combine before they can
be used. That notch is an interaction, not decoration on the consumer, so it is
modelled here as its own node --- \texttt{ix\_industry} and
\texttt{ix\_consumer}. Making it a node rather than an ornament is what lets the
model say which flows combine, which is the semantic content the notch carries.

\emph{Exchanges are modelled as single transactions, not as two pathways.} In
the original, money and energy meet at exchange diamonds: a purchase is one
transaction with an energy side and a money side. Declaring those as independent
pathways that happen to pass near a diamond would render similarly and assert
something weaker --- that a payment and a delivery both occurred, rather than
that they are the same event. A transaction node here names a seller and a
buyer, and the four half-pathways are generated from it. Note the design
consequence worth stating: because generation happens before validation, the
synthesised legs are checked and routed exactly like declared ones, so the
construct adds expressiveness without adding a special case anywhere downstream.

\begin{figure}[htbp]
\lstinputlisting[linerange={7-36}]{../examples/aggregated-economy-native.tex}
\caption{The aggregated economy authored in \TeX{}
(\texttt{examples/aggregated-economy-native.tex}).}
\label{fig:complex-tex}
\end{figure}

\begin{figure}[htbp]
\lstinputlisting[basicstyle=\ttfamily\scriptsize]{../examples/aggregated-economy.json}
\caption{The same model as JSON (\texttt{examples/aggregated-economy.json}).
Key names are identical to the \TeX{} form; only the syntax differs.}
\label{fig:complex-json}
\end{figure}

\begin{figure}[htbp]
\centering
\begin{subfigure}{\textwidth}
\centering
\includegraphics[width=0.86\textwidth]{aggregated-economy.png}
\caption{Odum's published figure.}
\end{subfigure}

\vspace{1em}
\begin{subfigure}{\textwidth}
\centering
\includegraphics[width=0.86\textwidth]{diagram-aggregated-economy.png}
\caption{Rendered by the package from either listing above.}
\end{subfigure}
\caption{A complete diagram beside its source. The system window frames what the
model treats as inside the system; dissipation pathways converge on a single
ground symbol placed outside it, because used energy leaves the system.}
\label{fig:complex-compare}
\end{figure}

\subsection{Division of responsibility}
\label{sec:raci}

[OUTLINE] State plainly where each obligation sits, because a reader deciding
whether to trust a figure needs to know what the renderer did and did not
check. The blunt version, which belongs in the paper verbatim: \emph{if the
model is energetically invalid, so is the diagram.} The package renders
faithfully; it does not repair.

\begin{table}[htbp]
\centering
\begin{tabular}{lll}
\toprule
Concern & Answerable & Notes \\
\midrule
Symbol geometry and conformance  & this package & measured against Odum's sheet \\
Placement conventions            & this package & \S\ref{sec:informal} \\
Topological energetic validity   & this package & four checks, \S\ref{sec:validity} \\
Model schema and identity        & shared       & one format, two implementations \\
Quantities and conservation      & GSSK         & the diagram carries no flux \\
Dynamics and simulation          & GSSK         & \\
Emergy accounting                & neither      & no agreed algorithm exists \\
Scientific validity of the model & the author   & no tool can supply this \\
\bottomrule
\end{tabular}
\caption{Where each obligation sits. The rows the package does not own matter as
much as those it does: a renderer that appeared to validate quantities would be
worse than one that plainly does not.}
\label{tab:raci}
\end{table}

Argue the last row rather than leaving it implicit. Odum's complaint was about
scholars, not software, and no amount of checking makes a model true. What
tooling can do is refuse to let one class of error pass silently, and be
explicit about which class that is.

\section{Using the package}
\label{sec:usage}

\subsection{Installation and invocation}
[OUTLINE] CTAN installation; the three authoring paths; a minimal complete
document.

\subsection{Model reference}
[OUTLINE] Node keys, edge keys, metadata keys, condensed to a table. Note that
key names are identical in both front ends by design.

\subsection{Naming and labelling}
[OUTLINE] Odum names symbols as well as describing them --- storages $Q_1,
Q_2, \ldots$, sources and flows $X$ and $J$ --- and both names carry
information: the short one is what the accompanying equations refer to, the long
one is what a reader needs. The package holds both, and
\texttt{label\_mode} selects which goes inside the symbol.

Make the argument for automatic displacement: a symbol sizes itself to its
label, but a label larger than the symbol's grid cell would stop it reading as
the cell it occupies --- the very inconsistency the grid exists to remove. Such
a label is therefore placed outside the symbol and, because it is then part of
what the diagram occupies, joins the obstacle set the router avoids.

\subsection{Ports, and the escape hatch}
[OUTLINE] The symbols are ordinary Ti\emph{k}Z shapes with documented
connection points, so a diagram that wants no automatic layout can be drawn by
hand and still use the correct symbols. Publishing those points is what makes
the package composable.

A symbol's connection points are a ring spaced a constant distance apart along
its outline (Figure~\ref{fig:anchors}), and they are its \emph{ports}: a
pathway attaches to one of them and to nothing else. The named anchors ---
\texttt{energy\_in}, \texttt{money\_out}, \texttt{heat\_sink} --- name a
\emph{direction} rather than a second, independent set of points: they say where
a kind of flow belongs, because that is where Odum draws it, and the pathway
attaches at the port nearest that direction. Making the two sets one is what
lets the renderer say anything about occupancy at all.

Argue why the ring is spaced by distance rather than by angle. Equal angular
divisions would put connection points further apart on a large symbol than on a
small one, so a coupled pair --- goods and money running together into the same
component --- would separate by a different amount at each symbol it met, and
the coupling would read inconsistently across one diagram. Constant spacing
makes the gap a property of the diagram rather than of the symbol, at the cost
of a variable number of ports per symbol, which is the right trade. The count is
rounded to a multiple of four so that a port falls on each axis, where Odum's
flows meet a symbol; that costs up to a quarter of the pitch in spacing and buys
an \texttt{energy\_in} that lands on the notch of an interaction rather than
half a pitch up its edge.

[OUTLINE] Then the occupancy rule, which is a claim about the diagram and not
about the drawing: \textbf{a port holds one pathway}. Two flows attached at the
same point are drawn as one line leaving the symbol, and the reader cannot
recover how many there were or which the arrowhead belongs to --- the diagram
asserts less than the model contains. So the second pathway to want a port takes
the next one along, on the side it approaches from. Note what this costs: the
attachment point is no longer a property of the symbol alone but of the symbol
and everything attached to it, so the allocation has to be independent of the
order the model happens to list its edges in --- our own two front ends list the
same diagram's exchanges in a different order, and \texttt{make parity} caught
it. Serving the pathways in the order of the directions they run off in, from
the middle of each fan outwards, is both order-independent and the condition for
a fan not to cross itself on the way out.

\begin{figure}[htbp]
\centering
\includegraphics[width=0.78\textwidth]{contact-sheet.png}
\caption{The implemented symbols with their ports: points spaced a constant
distance apart \emph{along the outline}, so a large symbol carries more of them
than a small one and a coupled pair attaching at adjacent indices is separated
by the same distance whichever symbol it meets. A pathway attaches to one of
these and to nothing else, and to one of them at a time; a named anchor such as
\texttt{energy\_in} resolves to the port nearest the direction it names. The
figure is rendered by the test suite, so a port that drifts off its outline is
caught --- which is how the switch's were found to sit up to a fifth of the
radius outside the bowed sides.}
\label{fig:anchors}
\end{figure}

\section{Interoperability with GSSK}
\label{sec:gssk}

[OUTLINE] Argue the payoff: a diagram and a simulation should not be two
independent descriptions of one system, because they will diverge.

GSSK (General Systems Simulation Kernel) takes a JSON model and integrates it.
Its schema and this package's arrived at nearly the same shape independently:
both use \texttt{metadata}, \texttt{nodes} and \texttt{edges}; and GSSK's node
type vocabulary --- \texttt{source}, \texttt{storage}, \texttt{sink},
\texttt{constant}, \texttt{interaction}, \texttt{gain}, \texttt{loop\_limited},
\texttt{exchange}, \texttt{switch}, \texttt{producer}, \texttt{consumer},
\texttt{misc\_box}, \texttt{system\_frame} --- is Odum's symbol set. The
convergence is not coincidence: both are formalising the same language.

Describe the alignment (\S\ref{sec:gssk} of the design note): energese reads
GSSK models directly, taking topology from the shared fields and presentation
from GSSK's \texttt{visual} object, which its schema already reserves and its
kernel ignores. The result is a single file that both simulates and draws, so a
published figure is generated from the model that produced the results beside
it.

\subsection{Round-tripping a visual editor}
\label{sec:roundtrip}

[OUTLINE] Make the case that automatic layout must be optional, not mandatory.
A modeller working in a drag-and-drop editor such as GSSK-Dia has already made
every placement decision, and a renderer that re-derives them publishes
something other than what was drawn.

The package therefore accepts geometry from the model at both levels: explicit
\texttt{x} and \texttt{y} on a node override the layout engine per axis, and
\texttt{waypoints} on an edge override the router. An editor can emit a model
that renders exactly as drawn, while a hand-written model carrying no geometry
gets the full layout treatment --- and the two mix freely, so a modeller may
place the components that matter and leave the rest to the engine.

State the consequence for the seam: the shared format then carries both the
simulation model and its visual presentation, which is what makes one file
sufficient for simulating and for publishing.

Adopting the shared vocabulary is what made the remaining symbol set explicit
rather than a matter of taste: GSSK names \texttt{gain}, \texttt{loop\_limited}
and \texttt{switch}, all three of which are now implemented and measured
(Table~\ref{tab:conformance}).

\section{Limitations and future work}
\label{sec:limits}

[OUTLINE] State plainly what remains, and be careful to claim no more than the
method supports.

\begin{itemize}
  \item \emph{Reference precision is the binding constraint.} Every tolerance in
        Table~\ref{tab:conformance} is set by the $144$\,dpi captures available,
        not by the renderer. The exchange diamond spans $31$ pixels, and at that
        size its own stroke consumes much of its interior. Higher-resolution
        sources would tighten every row, and are the single most valuable
        addition to this work.
  \item \emph{Conformance asserts two invariants, not a complete
        specification.} Aspect ratio and interior overlap together constrain a
        symbol tightly, but a shape could in principle satisfy both while being
        wrong in some detail neither captures. The contact sheet remains the
        human backstop.
  \item \emph{Some symbols are poor containers for text.} A triangle in
        particular: fitting the full bounding box of a label would force the
        amplifier to three times a producer's height, so the package guarantees
        the label's middle band and lets the extreme corners of the box fall
        outside. Odum labels amplifiers with a single character; long labels
        there will look wrong.
  \item \emph{Automatic layout handles the common cases.} Dense diagrams with
        many crossing feedback pathways still need per-edge routing, as
        Figure~\ref{fig:complex-tex} shows.
  \item \emph{Golden-image regression is engine-bound.} Goldens are only
        comparable across builds of the same \TeX{} Live, so an upgrade is a
        deliberate regeneration event rather than a transparent one.
\end{itemize}

\section{Toward a formal definition}
\label{sec:future-formal}

[OUTLINE] Where the paper says what should happen next; ambitious but specific.

\emph{A stated syntax.} The conventions implemented here are a candidate, not a
standard. What is missing is a specification the community has agreed, against
which independent implementations could be compared --- a position ESL has never
been in. Propose the checks of \S\ref{sec:validity} and the placement rules of
\S\ref{sec:informal} as a first draft to be argued with, and note that the
semiotic diagnosis of \S\ref{sec:semiotics} says where the work is needed:
syntactics, not semantics.

\emph{Emergy, deferred rather than dismissed.} Emergy --- embodied energy,
energy memory --- is central to the tradition and absent here. The reason is not
disagreement with the concept but the absence of an agreed algorithm, and
validating against a standard that does not exist would be worse than declining
to validate. Separate the two problems: depiction is tractable independently of
accounting. Valyi's emergy simulator showed dual energy and emergy flows on one
diagram,\footnote{% [PLACEHOLDER] Citation for Valyi's emergy simulator (Valyi
% and Ortega?). The notes describe a dual energy/emergy flow depiction; the
% bibliographic details need confirming.
Reference to be confirmed.} which is a notational proposal that could be adopted
without settling how emergy is computed. That is the next tractable step.

\emph{Quantities.} Conservation checking needs flux, which a diagram does not
carry. The shared model format (\S\ref{sec:gssk}) is the route, and it makes
this a question of interoperation rather than of scope creep.

\emph{Barrier to entry.} Return to \S\ref{sec:tooling}. The distribution problem
was never solved by better drawing tools; it is solved by a package registry, a
free implementation, and a format other tools can read. Argue that this is why
reviving the notation is plausible now in a way it was not when the only route
was a stencil emailed between research groups.

\section{Conclusions}
\label{sec:conclusions}

[OUTLINE] Return to the opening claim. An ESL diagram is a specification, so it
should be generated from a model and checked against its notation's conventions.
The package makes both routine. The wider methodological point --- separating
exact regression from tolerance-bounded conformance, and being explicit about
which claim the available reference material can support --- applies to any work
that reproduces figures from the printed scientific literature.

\section*{Software availability}

[OUTLINE] CTAN package name, repository URL, licence, version, archived DOI.
Note that every figure in this paper is regenerated by the package's own test
suite.

\begin{thebibliography}{99}

%% --- The Energy Systems Language -----------------------------------------
\bibitem[Odum(1971)]{odum1971}
Odum, H.T., 1971. \emph{Environment, Power, and Society}. Wiley-Interscience,
New York.

\bibitem[Odum(1994)]{odum1994}
Odum, H.T., 1994. \emph{Ecological and General Systems: An Introduction to
Systems Ecology}. University Press of Colorado, Niwot.

\bibitem[Odum(1996)]{odum1996}
Odum, H.T., 1996. \emph{Environmental Accounting: Emergy and Environmental
Decision Making}. Wiley, New York.

\bibitem[Brown and Ulgiati(2004)]{brown2004}
Brown, M.T., Ulgiati, S., 2004. Energy quality, emergy, and transformity:
H.T. Odum's contributions to quantifying and understanding systems.
\emph{Ecological Modelling} 178(1--2), 201--213.

%% --- Notation, semiotics, visual formalism --------------------------------
\bibitem[Morris(1938)]{morris1938}
Morris, C.W., 1938. \emph{Foundations of the Theory of Signs}. University of
Chicago Press, Chicago.

\bibitem[Goodman(1976)]{goodman1976}
Goodman, N., 1976. \emph{Languages of Art: An Approach to a Theory of Symbols},
2nd ed. Hackett, Indianapolis.

\bibitem[Bertin(1983)]{bertin1983}
Bertin, J., 1983. \emph{Semiology of Graphics: Diagrams, Networks, Maps}.
University of Wisconsin Press, Madison.

\bibitem[Harel(1988)]{harel1988}
Harel, D., 1988. On visual formalisms. \emph{Communications of the ACM} 31(5),
514--530.

\bibitem[Moody(2009)]{moody2009}
Moody, D.L., 2009. The ``physics'' of notations: toward a scientific basis for
constructing visual notations in software engineering. \emph{IEEE Transactions
on Software Engineering} 35(6), 756--779.

%% --- Implementation --------------------------------------------------------
\bibitem[Tantau(2023)]{tantau2023}
Tantau, T., 2023. \emph{The Ti\emph{k}Z and PGF Packages: Manual for Version
3.1.10}. Available from the Comprehensive TeX Archive Network.

%% -------------------------------------------------------------------------
%% [PLACEHOLDER] Still to be supplied by the author:
%%
%%   \bibitem{valyi}  Valyi's emergy simulator. The notes describe a tool that
%%       depicted dual energy and emergy flows on one diagram, and note that the
%%       formalisation of emergy it proposed was not widely adopted. Believed to
%%       be Valyi & Ortega; volume, year and venue need confirming before use.
%%
%%   \bibitem{hallTemplate}  The draw.io ESL template built at the request of
%%       Charles A. Hall for his doctoral students. Needs the repository URL and
%%       a date, plus whatever acknowledgement Hall prefers.
%%
%%   \bibitem{gssk}  GSSK (General Systems Simulation Kernel). Needs repository
%%       URL, version and, if one exists, a software citation or DOI.
%%
%%   Optionally: a reproducible-research citation supporting the
%%   figure-provenance argument in section 5, and citations for the
%%   systems-dynamics tools named in section 3 (Stella, Vensim, Insight Maker).
%% -------------------------------------------------------------------------

\end{thebibliography}

\end{document}
